% =============================================================================
% RQ2 (deobfuscation) section — REWRITTEN against the latest 2026-06-12 runs.
%
% Provenance: every number below is produced by
%   analysis/reproduce_rq2_deobfuscation.py
% from the pinned 6/12 CSVs:
%   DiffuCoder-7B_{all-masked,target-only}_20260612_102623.csv
%   DreamCoder-7B_{all-masked,target-only}_20260612_102631.csv
% and the RefineID targets in data/test.csv.
%
% WHAT CHANGED vs the V5 PDF (old -> new), all because data is now 6/12:
%   Table V  all-masked : DiffuCoder 12.2 -> 12.4 ; DreamCoder 13.9 -> 13.5
%   Table V  target-only: DiffuCoder  3.1 ->  2.8 ; DreamCoder  4.1 ->  4.7
%   Table VI : row id=9 (bytes) prediction  b -> f   (the only changed row; 8/9 identical)
%   Sec V-C  : 928 wrong -> 931 ; 96.9% -> 97.2% ; 662/71.3% -> 671/72.1% ;
%              189/20.4% -> 189/20.3% ; empty 77/8.3% -> 71/7.6% ;
%              example "parseInt" -> "randomGenerator" (parseInt no longer occurs; fConfig still does)
%   Sec V-D  : "3–4%" -> "2.8–4.7%" ; "71.3%" -> "72.1%" ;
%              "declines monotonically" -> corrected (the 1–10 bucket is BELOW the
%              11–20 bucket for DiffuCoder in BOTH the old and new data, so the
%              monotonic claim was never accurate — wording fixed here).
%   RQ1 clean ctx. (31.1 / 33.2) UNCHANGED — it comes from the separate RQ1
%              benchmark, not the deobfuscation runs.
%   Section III-D2 (metric definitions) UNCHANGED — methodology only, no numbers.
%   Abstract: "drops to about 3%" is still correct (target-only = 2.8% ≈ 3%); no
%              numeric edit needed. (The phrase "find the positions on their own"
%              is an abstract-wording matter, out of this RQ2 rewrite's scope.)
%
% Fig 8: regenerate from figures/new/fig8_rq2_em_june.{pdf,png}.
%   8(b) all-masked mean per-sample EM (%) by #identifiers bucket [1-10,11-20,21-50,51-100,100+]:
%     DiffuCoder-7B : 17.2  20.4  14.0   7.4   2.0
%     DreamCoder-7B : 21.5  19.9  13.7   7.5   2.4
% =============================================================================


% ---------- Table V ----------------------------------------------------------
\begin{table}[t]
  \centering
  \caption{Target-position Exact Match (\%) under RQ1 (clean context) and the
           two RQ2 experiments. Both RQ2 columns report EM evaluated at the
           \emph{same} \textsc{RefineID} target position, for direct
           comparability with RQ1.}
  \label{tab:rq2_em}
  \small
  \setlength{\tabcolsep}{6pt}
  \renewcommand{\arraystretch}{1.05}
  \begin{tabular}{@{}lccc@{}}
    \toprule
                  & RQ1            & \multicolumn{2}{c}{RQ2 target-EM} \\
    \cmidrule(l){3-4}
    Model         & clean ctx.     & all-masked & target-only \\
    \midrule
    DiffuCoder-7B & 31.1           & 12.4       & 2.8 \\
    DreamCoder-7B & 33.2           & 13.5       & 4.7 \\
    \bottomrule
  \end{tabular}
\end{table}


% ---------- Table VI ---------------------------------------------------------
\begin{table}[t]
  \centering
  \caption{DiffuCoder-7B's wrong predictions under the target-only condition.
           ``Obf.'' is the single-letter token assigned to the target during
           obfuscation. Rows marked $\dagger$ are cases where the model's
           prediction matches the obfuscated label of the target itself.}
  \label{tab:rq2_examples}
  \small
  \setlength{\tabcolsep}{6pt}
  \renewcommand{\arraystretch}{1.05}
  \begin{tabular}{@{}rlcc@{}}
    \toprule
    ID & Target & Obf. & Prediction \\
    \midrule
    0  & \texttt{decodedCapacity}   & ar & aq \\
    2  & \texttt{expectedIndexName} & n  & r  \\
    3  & \texttt{style}             & e  & e$^{\dagger}$ \\
    6  & \texttt{status}            & r  & a  \\
    9  & \texttt{bytes}             & j  & f  \\
    11 & \texttt{context}           & c  & a  \\
    12 & \texttt{cache}             & a  & a$^{\dagger}$ \\
    13 & \texttt{relativePath}      & j  & j$^{\dagger}$ \\
    15 & \texttt{printer}           & n  & p  \\
    \bottomrule
  \end{tabular}
\end{table}


% ---------- Section V prose --------------------------------------------------
\section{RQ2: Multi-Identifier Multi-Site Renaming}
\label{sec:rq2}

RQ1 showed that dLLMs can propagate a single renaming decision across many
co-referring sites in one bidirectional forward pass, and that the advantage
over FIM-AR baselines widens as the number of sites grows. The natural next
question is whether the same advantage extends from one identifier to many: can
a dLLM produce distinct names for several distinct identifiers in the same file
simultaneously? \textsc{RefineID} does not ask this question directly, since
each sample annotates exactly one identifier as the target. The experiment
setup in Section~III-D2 lets us decompose the gap between RQ1 and the full
multi-identifier multi-site task into two parts that can be measured separately.

\subsection{All-masked Experiment}
To probe renaming with multiple sites and multiple identifiers, we run the
all-masked experiment described in Section~III-D2. Under this condition,
DiffuCoder-7B reaches \textbf{12.4\%} target-EM and DreamCoder-7B reaches
\textbf{13.5\%} (Table~\ref{tab:rq2_em}). Both numbers are well below their RQ1
figures of 31.1\% and 33.2\%. Figure~8(b) reports how this number varies with
the number of distinct identifiers in the smell set $S$: per-sample EM is
highest for samples with only a handful of identifiers and falls off as $S$
grows, dropping to about 2\% for the most crowded files. This contrasts with
the RQ1 result on the same models (Figure~5), where the EM band of
26.8\%--36.4\% held across single-identifier multi-site samples regardless of
site count. The joint-resolution advantage is therefore robust to growth in the
number of co-referring sites of one identifier, but degrades as the number of
distinct identifier groups in $S$ grows.

\subsection{Target-only Experiment}
The second experiment keeps every identifier in place but rewrites each as a
single-letter token ($a, b, c, \dots$), so the model sees concrete
identifier-shaped tokens that carry no information about what those identifiers
mean, and only the target is masked. Under this condition, EM falls to
\textbf{2.8\%} (DiffuCoder-7B) and \textbf{4.7\%} (DreamCoder-7B)---9.6 and
8.8 percentage points below the all-masked result on the same target positions,
respectively. The two experiments differ only in whether the surrounding
identifiers are absent (all-masked) or actively misleading (target-only), so the
gap between them measures the cost of adversarial neighbours over and above the
cost of absent neighbours.

\subsection{Qualitative Analysis of Wrong Predictions}
The roughly 28 percentage points lost between clean context and target-only
leave the question of what DiffuCoder-7B actually produces in the 97.2\% of
target-only samples it gets wrong. We therefore inspect the wrong predictions.
Among the \textbf{931} wrong predictions, \textbf{671 (72.1\%)} are
one-or-two-letter lowercase strings that share the lexical form of the
obfuscated single-letter tokens that surround them. The remainder split between
longer meaningful identifiers (189, 20.3\%, such as \texttt{randomGenerator} and
\texttt{fConfig}) and empty strings from which no identifier can be extracted
(71, 7.6\%). Nearly three-quarters of the failures are therefore not informed
guesses that happen to disagree with the developer's choice---they are
predictions in the wrong shape. The general renaming ability of the model fails
in the target-only condition whenever there is surrounding obfuscation.

In three of the nine rows ($\dagger$) the prediction coincides with the
single-letter token the obfuscator had assigned to the target itself. This is a
chance collision, not recovery: the mapping is one-to-one and the target's
positions are all \texttt{<|mask|>}, so the assigned letter never appears in the
input. The bias toward frequent single-letter tokens under an obfuscated context
produces a small number of collisions with the alphabet. This bias is the
qualitative companion of the second decomposition step: the model in the
target-only condition follows the lexical style of the obfuscated context rather
than the structural cues that survive obfuscation, so masking those tokens away
(all-masked) removes the anchor and improves the result.

\subsection{Analysis}
The multi-site advantage observed in RQ1 was conditional on the model having
semantically informative neighbours to attend to. When the neighbours carry the
developer's chosen names, bidirectional attention can propagate one renaming
decision across many sites in a single pass, and the advantage scales with the
number of sites. When the same model is given concrete-but-wrong neighbours
(target-only), the EM advantage all but disappears, falling to 2.8--4.7\%. When
the neighbours are absent rather than misleading (all-masked), target-EM stays
low but still above the target-only result while below RQ1. Concrete-but-wrong
neighbours are therefore more costly than absent ones: 72.1\% of wrong
predictions under target-only adopt the lexical form of the obfuscated
neighbours, so the obfuscated context anchors the prediction to its own shape
rather than letting the model use the structural cues that survive obfuscation.

Figure~8(b) shows how the all-masked target-EM varies with the number of
distinct identifiers in $S$. EM falls off as the identifier count grows---from
the high-teens to low-20s for samples with at most a few dozen identifiers down
to about 2\% for files with 100 or more---which contrasts with the RQ1 result
(Figure~5), where the same dLLMs held an EM band of 26.8\%--36.4\% across
single-identifier multi-site samples regardless of site count. The advantage is
robust to growth in the number of co-referring sites of one identifier, but
degrades as the number of distinct identifier groups grows.

Therefore RQ2 supports RQ1 rather than contradicting it. The joint multi-site
advantage holds when the surrounding identifier names are semantically
informative and when $S$ resolves into a small number of co-reference groups,
and weakens when either condition fails.


% ---------- Fig 8 caption (unchanged; redraw from fig8_rq2_em_june.pdf) -------
% \caption{RQ2 results. (a) Target-position Exact Match for DiffuCoder-7B and
%   DreamCoder-7B under RQ1 (clean context), the all-masked experiment, and the
%   target-only experiment. (b) Mean per-sample Exact Match across
%   identifier-count buckets in the all-masked experiment.}
